\chapter{Visualization}
\label{ch:visualization}

BioDYM provides interactive Plotly-based visualizations in the Jupyter notebook. All plots follow a consistent publication style and can be exported as PNG (150 or 300~DPI), PDF, or SVG.

\section{Sankey Diagram}

The Sankey diagram provides a system-level overview of all material flows. Interactive controls:
\begin{itemize}
  \item \textbf{Element selector} --- Switch between material, WC, DM, CC
  \item \textbf{Year slider} --- View flows for any year in the model range
  \item \textbf{Threshold} --- Hide flows below a minimum value
\end{itemize}

Flow widths are proportional to mass. Node colors follow the process type color scheme.

\section{Process Dynamics (3-Panel View)}

For each process, a three-panel plot shows:
\begin{enumerate}
  \item \textbf{Top panel}: Inflows over time
  \item \textbf{Middle panel}: Outflows over time
  \item \textbf{Bottom panel}: Stock change ($\Delta S$) over time
\end{enumerate}

Select the process and element via dropdown menus.

\section{Flow Dynamics}

Time series plots of individual flows. Useful for verifying that input data was loaded correctly and that calculated flows behave as expected.

\section{Stock Analysis}

\begin{itemize}
  \item \textbf{Stock bar chart} --- Stacked bar showing stock levels across all processes at a selected year
  \item \textbf{System stock composition} --- Element breakdown of total system stock over time
\end{itemize}

\section{DSM Stock Details}

For processes with DSM logic, a specialized plot shows:
\begin{itemize}
  \item Stock evolution stacked by category (e.g., construction wood types)
  \item Selectable by process and element
  \item Includes decaying initial stock contribution
\end{itemize}

See Chapter~\ref{ch:dsm} for details on DSM configuration.

\section{Composition Analysis}

Interactive composition explorer showing the element fractions (WC, DM, CC) for each flow at each time step. Useful for verifying that compositions are physically consistent.

\section{Exporting Plots}

Each interactive plot has an export button or can be exported programmatically:

\begin{lstlisting}[style=python]
fig.write_image("my_plot.png", width=1200, height=800, scale=2)
\end{lstlisting}

\begin{tipbox}
For publication-quality figures, use \texttt{scale=3} (300~DPI at 1200px width). The publication style module automatically applies consistent fonts, colors, and layout.
\end{tipbox}

\section{Color Scheme}

BioDYM uses a consistent color palette:

\begin{table}[h]
\centering
\begin{tabular}{lll}
\toprule
\textbf{Purpose} & \textbf{Color} & \textbf{Hex Code} \\
\midrule
Material element & Blue & \texttt{\#2E86AB} \\
WC element & Green & \texttt{\#28A745} \\
DM element & Gray & \texttt{\#6C757D} \\
CC element & Yellow & \texttt{\#FFC107} \\
Splitter process & Blue & \texttt{\#2E86AB} \\
Transformer process & Pink & \texttt{\#A23B72} \\
\bottomrule
\end{tabular}
\caption{Standard color palette.}
\end{table}
